\documentclass[11pt]{article}
\usepackage[a4paper,margin=1in]{geometry}
\usepackage{booktabs}
\usepackage{listings}
\usepackage{amsmath}
\usepackage{hyperref}

\title{Problem 2 Full Report: OpenMP Acceleration of NEP Potential Calculation}
\author{}
\date{June 20, 2026}

\begin{document}
\maketitle

\section{Objective and Requirements}
Problem 2 requires OpenMP acceleration of molecular-dynamics potential
calculation. The work includes parallel-opportunity analysis, hotspot
identification, correct treatment of shared data, performance testing with
different thread counts, serial/parallel correctness comparison, and unit
testing. This work focuses on the NEP CPU implementation in
\texttt{source/source\_md/potential/ml/nep}.

\section{Parallel Analysis}
The NEP calculation contains neighbor-list construction, descriptor and neural
network derivative calculation, and radial/angular/ZBL force accumulation.
The first two stages already contain atom-level OpenMP loops. The force stage
was the main unfinished hotspot. For center atom $n_1$ and neighbor $n_2$:
\[
F_{n_1}\mathrel{+}=f_{12},\qquad F_{n_2}\mathrel{-}=f_{12}.
\]
Virial is also accumulated by atom index. A direct parallel loop therefore has
write/write conflicts when threads update the same neighbor.

\section{Three NEP Force Versions}
The NEP API now exposes \texttt{NEP::ForceParallelMode} and
\texttt{set\_force\_parallel\_mode}. The default remains
\texttt{ThreadLocal}, preserving optimized behavior for existing callers.

\subsection{Version 1: Serial Force Baseline}
\texttt{ForceParallelMode::Serial} executes radial, angular, and ZBL force
loops serially. It provides the deterministic reference. Neighbor and
descriptor stages may still use OpenMP, so additional threads can produce
limited speedup outside force accumulation.

\subsection{Version 2: OpenMP Atomic Accumulation}
\texttt{ForceParallelMode::Atomic} parallelizes the outer atom loop and uses
\texttt{omp atomic update} for every shared force, virial, and ZBL-energy
addition. It uses little extra memory, but synchronization occurs repeatedly
inside the neighbor loop.

\subsection{Version 3: OpenMP Thread-Local Reduction}
\texttt{ForceParallelMode::ThreadLocal} gives each thread private force and
virial arrays; ZBL also receives private energy storage. Threads accumulate
without synchronization, then reduce their buffers once after the loop. Its
memory cost is $O(TN)$, but it avoids atomics in the hotspot.

All versions share the same force mathematics; only accumulation changes. The
radial, angular, and ZBL small-box force kernels support all three modes.

\section{Correctness Testing}
The test \texttt{source/source\_md/test/nep\_omp\_test.cpp} loads
\texttt{tests/PP\_ORB/nep\_hfo2.txt}. It evaluates the same Hf/O system with
serial, four-thread atomic, and four-thread thread-local modes. Potential,
force, and virial use tolerances $10^{-12}$, $10^{-10}$, and $10^{-10}$.

The sweep also compares every result with one-thread serial output. Across all
27 cases:
\[
\max|\Delta E|=0,\quad
\max|\Delta F|=1.83187\times10^{-15},\quad
\max|\Delta V|=3.55271\times10^{-14}.
\]
These are floating-point summation-order differences and are well below the
test tolerances.

\section{Benchmark Script}
The benchmark consists of:
\begin{lstlisting}
source/source_md/tools/nep_omp_sweep.cpp
source/source_md/tools/run_problem2_nep_sweep.sh
\end{lstlisting}
It generates deterministic alternating Hf/O cubic systems, compiles with
\texttt{-O2 -fopenmp}, runs all modes with 1, 2, and 4 threads, checks each
result, and writes CSV. The command used was:
\begin{lstlisting}[language=bash]
cd source/source_md
tools/run_problem2_nep_sweep.sh ../../tests/PP_ORB/nep_hfo2.txt \
  64,216,512 1,2,4 20 reports/problem2_nep_sweep_results.csv
\end{lstlisting}
Model loading is excluded from timing. Full raw results are stored in
\texttt{reports/problem2\_nep\_sweep\_results.csv}.

\section{Performance Results}
Speedup is relative to serial-force, one-thread execution for each size.

\begin{table}[ht]
\centering
\small
\begin{tabular}{@{}llrrrr@{}}
\toprule
Atoms & Mode & 1 thread (ms) & 2 threads (ms) & 4 threads (ms) & 4-thread speedup \\
\midrule
64  & Serial       & 3.3298  & 2.7041  & 2.2777  & 1.46 \\
64  & Atomic       & 3.3062  & 2.5729  & 2.0893  & 1.59 \\
64  & Thread-local & 3.2625  & 1.8249  & 1.2259  & 2.72 \\
\midrule
216 & Serial       & 10.3443 & 8.2511  & 7.3411  & 1.41 \\
216 & Atomic       & 10.5334 & 6.7192  & 5.0961  & 2.03 \\
216 & Thread-local & 10.1750 & 5.4566  & 2.8917  & 3.58 \\
\midrule
512 & Serial       & 30.1092 & 23.1361 & 20.4003 & 1.48 \\
512 & Atomic       & 44.0341 & 19.8320 & 16.7979 & 1.79 \\
512 & Thread-local & 27.4264 & 16.0081 & 10.0029 & 3.01 \\
\bottomrule
\end{tabular}
\caption{Comparison of three NEP force-accumulation versions.}
\end{table}

\section{Analysis}
The serial-force version scales modestly because neighbor and descriptor work
is parallel while force remains a bottleneck. Its four-thread speedup stays
below $1.5\times$. Atomic mode improves scaling but shared-array
synchronization limits it; its 512-atom one-thread result also shows atomic
instruction overhead without contention.

Thread-local reduction is fastest overall. At four threads it reaches
$2.72\times$, $3.58\times$, and $3.01\times$ speedup for 64, 216, and 512
atoms. Moving synchronization outside the neighbor loop is therefore more
effective than atomic updates for every pair contribution.

\section{Requirement Coverage and Conclusion}
This implementation identifies dependencies and hotspots; parallelizes radial,
angular, and ZBL force calculations; provides serial, atomic, and thread-local
versions; handles races through atomics or private reduction; verifies numerical
consistency; benchmarks three sizes and thread counts; and supplies a repeatable
script and CSV.

Thread-local reduction is recommended as the default. Atomic mode provides a
low-memory comparison, and serial mode remains the correctness baseline. Future
work may reduce private-buffer memory using atom blocking or domain
decomposition and test scaling beyond four cores.

\end{document}
